\section{Greedy Function Inclusions}

Having a current set of formulas $\formulaset$ we want to choose the best $\formula\in\greedyhypothesis$ to extend the set of formulas to $\formulaset\cup\{\formula\}$ in a way minimizing the cross entropy.
Given this, add each step we solve the greedy cross entropy minimization
\begin{align}
    \label{prob:perfectGreedy}\tag{$\probtagtypeinst{\greedysymbol}{\datamap,\hlnstat,\greedyhypothesis}$}
    \min_{\formula\in\greedyhypothesis} \min_{\hybridparam\in\hybridparamsetofdim{\cardof{\hlnstat}+1}}
    \centropyof{\empdistribution}{\probof{\hlnstat\cup\{\formula\},\hybridparam}} \, .
\end{align}
A brute force solution of \probref{prob:perfectGreedy} would require parameter estimation for each candidate in $\greedyhypothesis$.
We provide two more efficient approximative heuristics in the following (see Chapter~20 in \cite{koller_probabilistic_2009}), which are faster to compute estimates of the cross entropy improvement from adding a formula to an existing statistic.

\subsection{Gain Heuristic}

In the gain heuristic, only the parameters of the new formula are optimized and the others left unchanged.
Let $\sechybridparam$ be the canonical parameter of the reference distribution on the statistic $\hlnstat$.
When adding a feature $\formula\in\greedyhypothesis$ we extend $\hlnstat$ by $\formula$ and restrict the parameters to coincide with $\sechybridparam$ on the first $\cardof{\hlnstat}$ coordinates.
We denote this constraint by
\begin{align*}
    \restrictionofto{\hybridparam}{[\cardof{\hlnstat}]} = \sechybridparam \, .
\end{align*}
The greedy gain heuristic is then the problem
\begin{align}
    \label{prob:greedyGain}\tag{$\probtagtypeinst{\gainsymbol}{\datamap,\formulaset,\sechybridparam,\greedyhypothesis}$}
    \min_{\formula\in\greedyhypothesis} \,\, \min_{\hybridparamsetofdim{\cardof{\hlnstat}+1} \wcols \restrictionofto{\hybridparam}{[\cardof{\hlnstat}]} = \sechybridparam }
    \centropyof{\empdistribution}{\probof{\hybridparam}} \, .
\end{align}
%Here we denote by $\canparam$ the first $\cardof{\formulaset}$ coordinates of the M-projection $\currentdistribution$  of $\empdistribution$ onto $\formulaset$ and the variable new coordinate at position $\canparamat{\cardof{\formulaset}}$.
To provide further insight into the solution of the gain heuristic, let us quantify the improvement of the cross entropy when adding a feature $\exformula$.

%% FOR GAIN HEURISTIC
\begin{lemma}
    Let $\empdistribution$, $\secprobtensor$ be distributions and $\exformula$ a formula such that $\contraction{\secprobtensor,\exformula}\notin\{0,1\}$.
    Then we have
    \begin{align*}
        \centropyof{\empdistribution}{\probof{\hlnstat,\sechybridparam}}
        - \min_{\acttensor} \centropyof{\empdistribution}{\probof{\acttensor}}
        = \kldivof{\bernoulliof{\contraction{\empdistribution,\exformula}}}{\bernoulliof{\contraction{\secprobtensor,\exformula}}}
    \end{align*}
    where by $\bernoulliof{p}$ we denote the Bernoulli distribution with parameter $p\in[0,1]$.
\end{lemma}
\begin{proof}
    We use the characterization of the local update by \lemref{lem:localHybridParamUpdate} for $\secprobtensor=\probof{\hlnstat,\sechybridparam}$.
    The update on the added feature is then parametrized by the two-dimensional vector $\acttensorat{\formulavar}$ and we have
    \begin{align*}
        \probof{\hlnstat,\sechybridparam} = \probof{\ones} \, .
    \end{align*}
    Let us abbreviate $\datamean\coloneqq\contraction{\empdistribution,\exformula}$ and $\currentmean\coloneqq\contraction{\probof{\hlnstat,\sechybridparam},\exformula}$.
    We distinguish the cases $\datamean\in(0,1)$ and $\datamean\in\{0,1\}$.

    In the case $\datamean\in(0,1)$, we have $\hardlegset=\sechardlegset$, $\headindexof{\hardlegset}=\secheadindexof{\hardlegset}$ and $\canparamat{\indexedselvariable}=\seccanparamat{\indexedselvariable}$ for $\selindex\in[\cardof{\hlnstat}]$.
    The additional coordinate of the canonical parameter is then
    \begin{align*}
        \canparamat{\selvariable=\cardof{\hlnstat}} = \lnof{
            \frac{\datamean}{(1-\datamean)}
            \cdot \frac{1-\currentmean}{\currentmean}
        } \, .
    \end{align*}
    Based on this characterization, the cross entropy difference is
    \begin{align*}
        \centropyof{\empdistribution}{\probof{\ones}} - \min_{\acttensor} \centropyof{\empdistribution}{\probof{\acttensor}} \\
        &= \datamean \cdot \canparamat{\selvariable=\cardof{\hlnstat}}
        - \lnof{\contraction{\secprobtensor,\formulaccwith,\softactsymbolofat{\cardof{\hlnstat},\canparam}{\formulavar}}} \, .
    \end{align*}
    We simplify
    \begin{align*}
        \contraction{\secprobwith,\formulaccwith,\softactsymbolofat{\cardof{\hlnstat},\canparam}{\formulavar}}
        = (1-\currentmean) + \currentmean \cdot \expof{\canparamat{\selvariable=\cardof{\hlnstat}}}
        = \frac{(1-\currentmean)}{(1-\datamean)} \, .
    \end{align*}
    We further have
    \begin{align*}
        \datamean \cdot \canparamat{\selvariable=\cardof{\hlnstat}}
        = \datamean \cdot \left[ \lnof{\frac{\datamean}{(1-\datamean)}\cdot \frac{(1-\currentmean)}{\currentmean}}  \right]
        = \datamean \lnof{\datamean} - \datamean \lnof{1-\datamean} + \datamean \lnof{1-\currentmean} - \datamean \lnof{\currentmean}
    \end{align*}
    and arrive at
    \begin{align*}
        &  \centropyof{\empdistribution}{\probof{\ones}} - \min_{\acttensor} \centropyof{\empdistribution}{\probof{\acttensor}} \\
        & \quad =  \datamean \lnof{\datamean} - \datamean \lnof{1-\datamean} + \datamean \lnof{1-\currentmean} - \datamean \lnof{\currentmean}
        -  \lnof{1-\currentmean} - \lnof{1-\datamean} \\
        & \quad = \left( -\datamean \lnof{\currentmean} - (1-\datamean) \lnof{1-\currentmean} \right)  - \left( -\datamean \lnof{\datamean} - (1-\datamean) \lnof{1-\datamean} \right) \\
        & \quad = \kldivof{\bernoulliof{\datamean}}{\bernoulliof{\currentmean}} \, .
    \end{align*}

    In the case $\datamean\in\{0,1\}$, the optimal $\acttensor$ is boolean and only the partition function term is changed by the update.
    We then have
    \begin{align*}
        \centropyof{\empdistribution}{\probof{\ones}} - \min_{\acttensor} \centropyof{\empdistribution}{\probof{\acttensor}}
        %&= \lnof{\contraction{\secprobwith}} - \lnof{\contraction{\secprobwith,\formulaccwith,\acttensorat{\formulavar}}} \\
        &= - \lnof{\contraction{\secprobwith,\formulaccwith,\acttensorat{\formulavar}}} \\
        &= \begin{cases}
               \lnof{\currentmean} & \ifspace \datamean=0 \\
               \lnof{1-\currentmean} & \ifspace \datamean=1 \\
        \end{cases} \\
        &= \kldivof{\bernoulliof{\datamean}}{\bernoulliof{\currentmean}} \, . \qedhere
    \end{align*}
\end{proof}

Problem \eqref{prob:greedyGain} is thus solved for
\begin{align*}
    \hat{\formula} \in \argmax_{\formula\in\formulaset} \kldivof{\bernoulliof{\contraction{\empdistribution,\formula}}}{\bernoulliof{\contraction{\currentdistribution,\formula}}}
\end{align*}
and $\hybridparam$ characterized in \lemref{lem:localHybridParamUpdate}.
The minimum is taken at
\begin{align*}
    \centropyof{\empdistribution}{\probof{\hlnstat,\sechybridparam}}
    - \kldivof{\bernoulliof{\contraction{\empdistribution,\hat{\formula}}}}{\bernoulliof{\contraction{\currentdistribution,\hat{\formula}}}} \, .
\end{align*}

% Algorithmic
The gain heuristic thus searches for the mode of a coordinatewise transform of the mean parameter tensors to $\empdistribution$ and $\currentdistribution$, using the Bernoulli Kullback-Leibler divergence as transform function.

% Interpretation
One therefore takes the formula, which marginal distribution in the current model and the targeted distribution are differing at most, measured in the KL divergence.

% Optimization method
%One optimization method would thus be the computation of the mean parameters to both distribution, building the coordinatewise KL divergence and choosing the maximum.
%Since we need to evaluate each coordinate, this can be intractable for large sets of formulas.


% Further weight optimization
Further improvement of the model can be achieved by iteratively optimizing the other weights as well, since their corresponding moment matching conditions might be violated after the integration of a new formula.
Unfortunately, backward mappings cannot be expressed in closed form in general.
%This would require the computation of backward mappings for each candidate formula, for which we only have an alternating approach in general.



\subsection{Gradient Heuristic}

Gradient heuristic is another approach to select a feature.
We show here how the effective selection tensor network representation of exponentially many formulas described in \charef{cha:formulaSelection} can be utilized. % Or later in proposal

In the gradient heuristic, we estimate the cross entropy improvement by the gradient of the cross entropy with respect to varying with another feature.
Given a distribution $\secprobwith$, we vary % more general than
\begin{align*}
    \probofat{\softacttensor}{\shortcatvariables}
    \coloneqq \normalizationof{\secprobwith,\formulaccwith,\softacttensorat{\headvariableof{\formulavar}}}{\shortcatvariables} \, .
\end{align*}
In the gradient heuristic we choose the steepest gradient direction
\begin{align}
    \label{prob:greedyGrad}\tag{$\probtagtypeinst{\gradientsymbol}{\datamap,\secprobtensor,\greedyhypothesis}$}
    \min_{\formula\in\greedyhypothesis} \difofwrt{\centropyof{\empdistribution}{\probofat{\softacttensor}{\shortcatvariables}}}{\canparam}   \, .
\end{align}

%% FOR GRADIENT HEURISTIC
\begin{lemma}
    For any formula $\formula$ we have
    \begin{align*}
        \difofwrt{\centropyof{\empdistribution}{\probof{\softacttensor}}}{\canparam} =
        -\contraction{\empdistribution,\formula} + \contraction{\probof{\softacttensor},\formula} \, .
    \end{align*}
\end{lemma}
\begin{proof}
    We have
    \begin{align*}
        \centropyof{\empdistribution}{\probof{\softacttensor}}
        & = \contraction{\empdistribution,\formulaccwith,-\lnof{\softacttensorat{\headvariableof{\formulavar}}}}
        + \lnof{\contraction{\secprobwith,\formulaccwith,\softacttensorat{\headvariableof{\formulavar}}}}
    \end{align*}
    The derivation of the first term at $\canparam=0$ is $-\contraction{\empdistribution,\formula}$ and of the second $\contraction{\probof{\softacttensor},\formula}$.
\end{proof}

\probref{prob:greedyGrad} is thus
\begin{align}
    %\label{prob:greedyGrad}\tag{$\probtagtypeinst{\gradsymbol}{\datamap,\secprobtensor,\greedyhypothesis}$}
    \min_{\formula\in\greedyhypothesis} \contraction{\secprobtensor,\formula} - \contraction{\empdistribution,\formula}   \, .
\end{align}
%% Positive and Negative Search
The gradient shows the typical decomposition into a positive and a negative phase.
While the positive phase comes from the data term and prefers directions of large data support, the negative phase originates in the partition function and draws the gradient away from directions already supported by the current model $\expdistof{(\naivestat, \naivecanparam)}$.
%% Regularization functionality
The negative phase is a regularization, by comparing with what has already been learned.
When nothing has been learned so far, we can take the current model to be the uniform distribution, which is the naive exponential family with vanishing canonical parameters.






\subsection{Proposal Distributions}

We now frame the selection of a formula as a sampling problem of proposal distributions.
%All the costs and the exact greedy objective
Each of the scores to candidate formulas are understood as coordinates of an energy tensor.
Among the heuristics, the gradient heuristic has the most accessible form, since the energy tensor is available as a tensor network of the selection encoding of the formulas.

%We can choose selection architectures to efficiently parametrize the formulas in the hypothesis $\greedyhypothesis$ and rewrite the problem as
%\begin{align*}
%	\argmax_{\selindexin} \contractionof{ \gradwrtat{\canparam}{\canparam=0} \lossof{\expdist}}{\indexedselvariable}
%\end{align*}
%This is thus equivalent to the problem \ref{prob:greedyGrad}, when taking all formulas selectable by $\formulaset$ as the hypothesis $\Gamma$.

% Proposal distribution
%Let us now understand the likelihood gradient as the energy tensor of a probability distribution, which we call the proposal distribution.

\begin{definition}[Gradient Heuristic Proposal Distribution]
    Let there be a base distribution $\currentdistribution$, a targeted distribution $\empdistribution$ and a formula selecting map $\greedyhypothesis$.
    The proposal distribution at inverse temperature $\invtemp>0$ is the distribution of $\selvariable$ defined by
    \begin{align*}
        \normalizationof{\expof{\contractionof{\invtemp\cdot(\empdistributionwith-\currentdistributionwith),\greedyhypothesis\left[\shortcatvariables,\selvariable\right]}{\selvariable}} }{\selvariable} \, .
    \end{align*}
    The proposal distribution is the member of the exponential family with statistics $\greedyhypothesis$ and canonical parameter $\invtemp\cdot(\empdistribution-\currentdistribution)$.
\end{definition}


%. Exponential family
The proposal distribution is in the exponential family with sufficient statistic by the formula selecting map $\greedyhypothesis$, namely the member with the canonical parameters $\canparam=\empdistribution-\currentdistribution$.
Of further interest are tempered proposal distributions, which are in the same exponential family with canonical parameters $\invtemp\cdot(\empdistribution-\currentdistribution)$ where $\invtemp>0$ is the inverse temperature parameter.

% MLN
As \MarkovLogicNetworks{}, the proposal distributions are in exponential families with the sufficient statistic defined in terms of formula selecting maps.
While \MarkovLogicNetworks{} contract the maps on the selection variables $\selvariable$, the proposal distributions contract them along the categorical variables $\catvariable$ to define energy tensors.

% Methods to solve mode search
The gradient heuristic optimization \probref{prob:greedyGrad} is the search for the mode of the proposal distribution.
To solve the gradient heuristic, we thus need to answer a mode query, for which we can apply the methods introduced in \charef{cha:probReasoning}, such as Gibbs Sampling or Mean Field Approximations in combination with annealing.


%\subsection{Mean parameter polytope}
The mean parameter polytope of any proposal distribution to statistic $\proposalstat$ is the convex hull of the formulas in $\formulaset$, that is
\begin{align*}
    \meansetof{\proposalstat}
    = \convhullof{\sencodingof{\proposalstat}{\indexedselvariable,\shortcatvariables} \wcols \selindexin}
    = \convhullof{\formulaat{\shortcatvariables} \wcols \formula\in\greedyhypothesis}
\end{align*}
% 0/1
As it was the case for \MarkovLogicNetworks{}, the mean parameter polytopes are instances of a $0/1$-polytopes \cite{ziegler_lectures_2000,gillmann_01-polytopes_2007}.
% Interpretation as formulas
The vertices are the formulas selectable by the formula selecting map $\greedyhypothesis$.

\subsection{Iterations}

Let us now iterate the search for a best formula at a current model with the optimization of weights after each step.
The result is Algorithm~\ref{alg:greedyStructureLearning}, which is a greedy algorithm adding iteratively the currently best feature.

\begin{algorithm}[hbt!]
    \caption{Greedy Structure Learning}\label{alg:greedyStructureLearning}
    \begin{algorithmic}
        \Require Empirical distribution $\empdistribution$, hypothesis $\greedyhypothesis$ of formulas
        \Ensure Distribution $\expdist$ approximating $\empdistribution$
        \iosepline
        \State Initialize
        \[ \currentdistribution \algdefsymbol \frac{1}{\prod_{\catenumeratorin}\catdimof{\atomenumerator}} \cdot \onesat{\shortcatvariables} \quad, \quad \formulaset = \varnothing \]
        \While{Stopping criterion is not met}
        % REFINE! Work in data
            \State
            \begin{itemize}
                \item \textbf{Structure Learning:} Compute a (approximative) solution $\hat{\formula}$ to \probref{prob:perfectGreedy} and add the formula to $\formulaset$, i.e.
                \[ \formulaset \algdefsymbol \formulaset \cup\{\hat{\formula}\} \]
                Extend dimension of $\selvariable$ by one, by $\formulaof{\seldim}=\hat{\formula}$ and $\canparamat{\selindex=\seldim}=0$
                \item \textbf{Weight Estimation:} Estimate the best weights for the added formula and recalibrate the weights of the previous formulas, by calling Algorithm~\ref{alg:AMM_HLN}.
                \[ \currentdistribution \algdefsymbol \probof{\formulaset,\hybridparam} \]
            \end{itemize}
        \EndWhile
        \State \Return $\formulaset$, $\hybridparam$ %, $\kb$
    \end{algorithmic}
\end{algorithm}


%% Energy Storage -> Useful after learning for energy-based inference
When having used the same learning architecture multiple times, the energy of the corresponding formulas are all representable by a formula selecting architecture.
Their energy term is therefore a contraction of the selecting tensor with a parameter tensor $\canparam$ in a basis $\cpformat$ decomposition with rank by the number of learned formulas.
When multiple selection architectures have been used, the energy is a sum of such contractions.
% 
Let us note, that this representation is useful after learning, when performing energy-based inference algorithms on the result.
During learning, one needs to instantiate the proposal distribution, which requires instantiation of the probability tensor.
\red{However, one could alternate data energy-based and use this as a particle-based proxy for the probability tensor.}


\begin{remark}[Sparsification by Thresholding]
    To maintain a small set of active formulas, one could combine greedy learning approaches with thresholding on the coordinates of $\canparam$.
    This is a standard procedure in Iterative Hard Thresholding algorithms of Compressed Sensing, but note that here we do not have a linear in $\canparam$ objective.
\end{remark}


%\sect{Discussion}
%
%\begin{remark}[Bayesian approach]
%    We only treated the estimation of a single resulting distribution by the data, while in a Bayesian approach one typically considers an uncertainty over possible distributions.
%    % MAP
%    \red{When treating $\canparam$ as a random tensor, which prior distribution is given and posteriori distribution wanted, we have a more involved Bayesian approach.}
%    When having a prior $\probat{\mlnparameters}$ over the \MarkovLogicNetworks{} we alternatively want to find the parameters $\mlnparameters$ solving the maximum a posteriori problem
%    \begin{align}
%        \argmax_{\mlnparameters} \mlnprobat{\data}\cdot \probat{\mlnparameters}\, .
%    \end{align}
%\end{remark}
